\section{Methodology} \label{sec:methodology}
  % In General
  This chapter describes the methodological framework applied in this thesis 
  with the focus placed on theory, analysis and evaluation of the experiments. 
  It outlines the research design, the applied methodological approach, as 
  well as the procedures for data collection, analysis, and evaluation used 
  to assess the susceptibility of modern Raspberry Pi models to Rowhammer 
  attacks.

  \subsection{Research Design} \label{subsec:research-design}
    % Research Type & Objective
    This study is quantitative using experimental testing to confirm or refute 
    specific hypotheses in order to obtain empirical evidence and clarity 
    regarding the susceptibility of various modern Raspberry Pi models to 
    Rowhammer attacks.
    \\
    % Methodological Approach
    From a methodological perspective, the study primarily adopts a deductive approach. We tested existing 
    Rowhammer theories on Raspberry Pi platforms by starting from established theoretical 
    foundations, formulating hypotheses derived from these theories, and subsequently validating or 
    refuting them through controlled experiments. In general, this deductive approach is characterized by 
    the investigation of an existing theory, followed by the derivation of hypotheses, the execution of 
    experiments, and the confirmation or falsification of the proposed assumptions. \\
    In addition to the deductive core, the methodology incorporates an inductive approach in a supporting 
    role. Experimental observations obtained during testing may reveal unexpected system behavior. Based 
    on these measured results, new hypotheses or insights can be derived. In general terms, the inductive 
    approach involves the derivation of theory from empirical research by identifying patterns in 
    observations and formulating new hypotheses or theoretical considerations. Overall, we classify the study 
    predominantly deductive and complement it with inductive elements that emerge from experimental 
    findings. 
    \\
    % Test Devices
    The experimental setup included Raspberry Pi 4 models with 4~GB and 8~GB of RAM, as well as 
    Raspberry Pi 5 models with 4~GB, 8~GB, and 16~GB of RAM (see~\cref{tab:experimental-setups}). 
    \\
    The Raspberry Pi 4 devices are equipped 
    with a Quad Core ARMv8 Cortex-A72 \gls{cpu} at 1.8~GHz, featuring 32~KB data and 48~KB instruction L1 
    cache per core and a 1~MB L2 cache, which collectively provide approximately 50~\% higher performance 
    than the Raspberry Pi 3B~\cite{webProcessors}. 
    \\
    The Raspberry Pi 5 devices include a Quad Core ARMv8 Cortex-A76 \gls{cpu} 
    at 2.4~GHz, with 64~KB data and instruction L1 cache per core, 512~KB L2 cache per core, and a 
    shared 2~MB L3 cache. In aggregate, the new features present in BCM2712 deliver a performance 
    uplift of 2-3x over Raspberry Pi 4 for common \gls{cpu} or I/O-intensive use cases~\cite{webProcessors}.
    The selected ARM-based single-board computers were provided for the purpose of this study.
    The devices are comparable in terms of architecture, and overall system design, which 
    enables a controlled comparison of Rowhammer susceptibility.

    \begin{table}[h!]
      \centering
      \caption{Evaluated Raspberry Pi configurations}
      \label{tab:experimental-setups}
      \begin{tabular}{|l|l|l|l|l|}
        \hline
        \textbf{Platform} & \textbf{Processor} & \textbf{Architecture} & \textbf{SoC} & \textbf{DRAM} \\ \hline
        Raspberry Pi 4B & ARMv8 & aarch64 & Broadcom BCM2711 & 4~GB LPDDR4 \\ \hline
        Raspberry Pi 4B & ARMv8 & aarch64 & Broadcom BCM2711 & 8~GB LPDDR4 \\ \hline
        Raspberry Pi 5 & ARMv8 & aarch64 & Broadcom BCM2712 & 4~GB LPDDR4x \\ \hline
        Raspberry Pi 5 & ARMv8 & aarch64 & Broadcom BCM2712 & 8~GB LPDDR4x \\ \hline
        Raspberry Pi 5 & ARMv8 & aarch64 & Broadcom BCM2712 & 16~GB LPDDR4x \\ \hline
      \end{tabular}
    \end{table}

    % Data Sources
    The primary data sources consisted of related academic papers, 
    providing the theoretical foundation. Supplementary 
    technical documentation, datasheets, and glossaries to clarify hardware specifications 
    and terminology. Additionally, GitHub repositories containing Rowhammer 
    proof-of-concept implementations were utilized, adapted, and applied to 
    the target devices.
    \\
    % Data Collection
    Data collection involved the experimental execution of Rowhammer attacks on each 
    Raspberry Pi model, combined with the observation of system behavior. Outcomes 
    were systematically recorded and logged to facilitate further hypothesis formation. 
    Data collection was conducted exclusively on the experimental devices to derive 
    conclusions from the observed results and to support the generation of subsequent 
    hypotheses.

  \subsection{Evaluation} \label{subsec:evaluation}
    To ensure the reliability and validity of results, experiments were repeated under 
    controlled environmental conditions to minimize external interference. Observed 
    outcomes were analyzed using statistical tests to assess the significance and
    consistency of the findings.
    \\
    % Ethical Considerations
    The experiments were conducted on the provided devices in a controlled setup. 
    The study focuses on the responsible, academic investigation of Rowhammer 
    vulnerabilities, with results intended solely for theoretical analysis and 
    evaluation, avoiding any risk to other systems.
